\section{Step 51 Name the two output representations}\label{h263-quant-s01}
Reserve separate 128-byte arrays for levels and reconstructed coefficients in a
private frame. The public result will contain both, but no public stores should
occur while metadata or sample-domain validation can still fail. Preserve the
saved registers you use and restore the caller's stack on every return path.

Keep the result pointer separate from the current input cursor. Both point at
halfwords, which makes an accidental pointer substitution plausible. Name their
roles in comments and inspect their addresses before debugging arithmetic.

\paragraph{Observe} The result has 256 active bytes. Neither private array
is a picture buffer, and neither stores final unsigned pixels.
\Needspace{8\baselineskip}
\paragraph{Check} Which array is ready to feed the inverse transform?
\begin{itemize}
\item[A.] The reconstructed coefficient array.
\item[B.] The array of transmitted levels without inverse quantization.
\item[C.] The eight-bit DC code by itself.
\end{itemize}

\section{Step 52 Admit complete controls before samples}\label{h263-quant-s02}
Load each request field as a doubleword. Check $1\leq Q\leq31$ and both binary
selectors before using them. An unsigned comparison of $Q-1$ with 31 can perform
the range check, including rejection of zero and large values whose low word
looks valid. Check pointer alignment and non-wrapping extents before dereferencing
samples. Mapped storage and disjointness are caller obligations.

Only after admission should you read signed halfwords. Validate all values against
$[-2048,2047]$. You may compute private levels as you validate, because an invalid
late sample still leaves public output untouched. Do not publish those partial
results.

\paragraph{Observe} Quantizer zero is rejected before any division or sample
read. A bad final coefficient returns status two with all 256 result bytes intact.
\Needspace{8\baselineskip}
\paragraph{Check} Why retain the full quantizer field during admission?
\begin{itemize}
\item[A.] It contains a video frame rate.
\item[B.] A large invalid value can have the same low word as eight.
\item[C.] It permits division by zero.
\end{itemize}

\section{Step 53 Select ordinary levels from magnitudes}\label{h263-quant-s03}
For an ordinary coefficient, form its nonnegative magnitude in a general register.
The most negative legal value is $-2048$, so its magnitude is 2048 and fits easily
in a word or doubleword. Divide this magnitude by the validated positive $2Q$,
cap at 127, and restore the original sign. Keep a zero quotient as zero.

Computing a magnitude first makes the dead zone symmetric. Arithmetic-shifting a
negative coefficient is not an interchangeable replacement for division of its
magnitude: a shift rounds negative values downward. The policy requires truncation
of the nonnegative quotient before restoring the sign.

\paragraph{Observe} At $Q=8$, coefficients 15 and $-15$ both become zero;
16 becomes one and $-16$ becomes minus one.
\Needspace{8\baselineskip}
\paragraph{Check} What is the level selected for $C=-15$ at $Q=8$?
\begin{itemize}
\item[A.] $-1$.
\item[B.] 15.
\item[C.] Zero.
\end{itemize}

\section{Step 54 Encode INTRA DC separately}\label{h263-quant-s04}
When INTRA is selected and the coefficient index is zero, use the eight-unit DC
policy instead. Add four, divide downward by eight, and clamp the selected $d$
to 1 through 254. An arithmetic right shift by three gives the stated downward
rounding for the signed numerator. Replace selected $d=128$ with code 255 in
the level record. Keep the distinction between selected $d$ and stored code.

Do not increment the ordinary nonzero count for this entry. Continue with the
ordinary policy for indices 1 through 63. In INTER mode, index zero follows the
ordinary policy and participates in the count.

\paragraph{Observe} INTRA coefficient 1024 stores code 255. An INTRA zero
block still stores DC code one, with ordinary nonzero count zero.
\Needspace{8\baselineskip}
\paragraph{Check} What reconstructed coefficient does DC code 255 represent?
\begin{itemize}
\item[A.] 1024.
\item[B.] 255.
\item[C.] 2040.
\end{itemize}

\section{Step 55 Reconstruct scalar magnitudes and signs}\label{h263-quant-s05}
Read each ordinary level from the private record. For zero, store zero directly.
For a nonzero value, calculate $Q(2|L|+1)$, subtract one if $Q$ is even, then
restore the level's sign. Clip the signed result to $[-2048,2047]$ and store it
in the private reconstruction array.

For INTRA DC, translate code 255 back to selected $d=128$, then multiply $d$ by
eight. This result bypasses the ordinary formula. You may process a placeholder
at index zero in a common loop if you replace it correctly before publication;
the public record must contain the special result.

\paragraph{Observe} At $Q=8$, levels one and minus one reconstruct 23 and
minus 23. Level zero reconstructs zero, not seven.
\Needspace{8\baselineskip}
\paragraph{Check} When is the even-parameter subtraction applied?
\begin{itemize}
\item[A.] After adding a predicted pixel.
\item[B.] To the positive nonzero magnitude before restoring the sign.
\item[C.] To the final pointer.
\end{itemize}

\section{Step 56 Establish lane widths before vectorizing}\label{h263-quant-s06}
Keep level selection scalar initially. The scalar path already handles the
single DC exception and the positive division. The mode-one path can then
reconstruct eight ordinary levels at once in halfword lanes.

For a legal ordinary magnitude at most 127, the largest un-clipped positive
reconstruction is $31(2\cdot127+1)=7905$. It fits a signed halfword. Even a
temporary DC code 255 processed by the same vector loop produces magnitude at
most 15841, also within that width; its result must be replaced by the special
DC reconstruction. These bounds justify halfword multiplication here. They would
not justify halfword products in the earlier transform's second pass.

\paragraph{Observe} Every intermediate magnitude in this vector reconstruction
fits the chosen lane width before clipping. This is an arithmetic bound, not an
assumption that the multiplication saturates.
\Needspace{8\baselineskip}
\paragraph{Check} What permits halfword multiplication in this stage?
\begin{itemize}
\item[A.] Every video operation uses bytes.
\item[B.] MSA automatically widens every product.
\item[C.] The legal input ranges bound the complete intermediate below 32768.
\end{itemize}

\section{Step 57 Restore vector signs and exact zeros}\label{h263-quant-s07}
Derive an all-zero or all-one halfword sign mask by arithmetic-shifting the signed
level fifteen places. XOR with that mask and subtract the mask to obtain its
magnitude. Apply the nonzero reconstruction arithmetic, then use the same XOR
and subtraction to restore signs.

Construct an equality mask for original zero levels and clear those lanes after
the formula. Without that mask, every zero lane acquires $Q$ or $Q-1$. Clamp
signed lanes to the coefficient limits. Load and store exactly the eight groups
that cover each 128-byte private array; no tail or neighbour data is needed.

\paragraph{Observe} Mixed positive, negative and zero lanes agree with the
scalar reconstruction at every position, not just in their sum.
\Needspace{8\baselineskip}
\paragraph{Check} Why retain the original zero-level mask?
\begin{itemize}
\item[A.] The nonzero formula would otherwise turn zero into a nonzero value.
\item[B.] Zero levels use a different memory byte order.
\item[C.] It replaces the input pointer check.
\end{itemize}

\section{Step 58 Publish both completed arrays}\label{h263-quant-s08}
Replace the INTRA DC reconstruction after the common vector loop if necessary.
Then copy both private arrays into the caller's result and return the ordinary
nonzero count. The selected level at coefficient zero can change this count only
in INTER mode.

Use the supplied default through the ordinary caller. Stop at the return label
before its first instruction and inspect the two arrays as signed halfwords.
The DC code 255 is stored as a positive halfword value, so it remains legible
alongside the signed ordinary levels. Byte packing belongs to the later bit writer.

\paragraph{Observe} The default returns status zero, ordinary nonzero count
seven, first stored code 255 and first reconstructed coefficient 1024.
\Needspace{8\baselineskip}
\paragraph{Check} Which value is excluded from the INTRA nonzero count?
\begin{itemize}
\item[A.] Every negative AC level.
\item[B.] The separate INTRA DC code.
\item[C.] Every coefficient in a chroma block.
\end{itemize}

\section{Step 59 Exercise boundaries and equivalent schedules}\label{h263-quant-s09}
Run the public feedback against your working source:
\begin{verbatim}
python3 construction/h263-quantization-construction/grade.py \
    my-quantizer.asm --cli /path/to/mars-cli.jar \
    --report quantizer-results.json
\end{verbatim}
The public cases test three quantizer parameters, both block types, modes and
byte orders, several input patterns, invalid fields, domain failures and state
preservation. Input begins at an address aligned for a halfword but not a vector,
so a direct vector load from the caller needs its own alignment reasoning.
Staging levels privately avoids that input alignment issue in the taught path.

Try another iteration order or private frame. Results and access boundaries
must stay the same. A different schedule is acceptable; the final exercise
contract does not require the reference's labels or scratch registers.

\paragraph{Observe} Successful feedback checks all 128 output halfwords. It
also verifies that a late invalid coefficient leaves the complete result unchanged.
\Needspace{8\baselineskip}
\paragraph{Check} What does exact kernel output agreement establish?
\begin{itemize}
\item[A.] The whole stream is already conforming H.263.
\item[B.] Every complete encoder must choose these same forward levels.
\item[C.] The tested kernel follows this exercise's quantization policy.
\end{itemize}

\section{Step 60 Feed reconstruction from reconstructed coefficients}\label{h263-quant-s10}
The next connection is from the second output array to the inverse-transform
entry. For INTER, add each signed inverse residual to its predicted sample in a
wide signed value, then clip the sum to 0 through 255. For INTRA, clip the inverse
output directly to the sample range. Keep these final pixel operations outside
the quantizer: its reconstruction values are still transform coefficients.

Retain both arrays until coefficient coding and local reconstruction have consumed
them. Coding will read the selected levels; the local picture path will read
the reconstructed coefficients. Later prediction must use that locally reconstructed
picture, including the changes caused by quantization and final clipping.

\paragraph{Observe} A coefficient of 16 at $Q=8$ becomes reconstructed
coefficient 23. Feeding 16 to the inverse would bypass the very loss the decoder sees.
\Needspace{8\baselineskip}
\paragraph{Check} Which coefficient block must the local inverse consume?
\begin{itemize}
\item[A.] The inverse-quantized block derived from the transmitted levels.
\item[B.] The original forward block regardless of selected levels.
\item[C.] The preceding raw picture's bytes reinterpreted as halfwords.
\end{itemize}
